\section{\texorpdfstring{\code{send()}}{send()}}
\label{sec:send}

\begin{cppcode}
bool Spi::send(const Frame& frame)
{
    if (!isValid(frame))                              { return false; }
    if ((readReg(RegStatus) & StatusTxReady) == 0U)   { return false; }

    writeReg(RegTxId,      frame.id);
    writeReg(RegTxDlc,     frame.dlc);
    writeReg(RegTxDataLo,  packDataLo(frame));
    writeReg(RegTxDataHi,  packDataHi(frame));
    writeReg(RegTxSend,    Trigger);

    return true;
}
\end{cppcode}

\subsection{The order is not arbitrary}

\code{TX_SEND} \textbf{must} be written last. That write does not mean "store a value" but "start
transmitting now", and the controller then reads \code{TX_ID}, \code{TX_DLC} and the two data
registers. Write the trigger first and you send the previous frame's contents --- which at
bring-up looks like a node running exactly one frame behind.

The status check has to come \textbf{before} the writes. If you were to write \code{TX_ID} while a
transmission is in progress, you would overwrite data the controller is in the middle of using.

\subsection{Two \texorpdfstring{\code{false}}{false}s, two meanings}

The interface has only one \code{bool} to express both with:

\begin{narrowtable}{@{}ll@{}}
\thead{Cause} & \thead{What the caller ought to do} \\ \midrule
Invalid frame & Fix the frame. Waiting never helps. \\
Hardware busy & Wait and try again. \\
\end{narrowtable}

That is a deliberate simplification, but one worth being aware of. Two things follow:

\textbf{The order matters.} The validation first, so that an invalid frame is rejected
\emph{without} a single transaction going out on the bus. The supplied test suite checks exactly
that.

\textbf{It becomes a problem exactly one seam further up.} This loop is an infinite loop for a
frame with \code{dlc = 9}:

\begin{cppcode}
while (!driver.send(frame))   // hangs forever if the frame is invalid
{
}
\end{cppcode}

The right handling is to validate before entering the loop, or to count attempts and give up.

\section{The packing}
\label{sec:packing}

Data byte 0 sits in the \textbf{most significant} bits, and \code{TX_DATA_LO} carries the
\textbf{first} four bytes.

\begin{cppcode}
/**
 * @brief Pack four payload bytes into one register word, byte 0 in the most significant bits.
 */
[[nodiscard]] std::uint32_t pack(const std::uint8_t* const data,
                                 const std::uint8_t offset) noexcept
{
    return (static_cast<std::uint32_t>(data[offset + 0U]) << 24U)
         | (static_cast<std::uint32_t>(data[offset + 1U]) << 16U)
         | (static_cast<std::uint32_t>(data[offset + 2U]) << 8U)
         | (static_cast<std::uint32_t>(data[offset + 3U]));
}
\end{cppcode}

The same cast rule as in \kapref{ch:byte}: \textbf{cast before the shift}.

The two most common packing errors:

\begin{booktable}{@{}lL@{}}
\thead{Error} & \thead{Symptom} \\ \midrule
\code{TX_DATA_LO} and \code{TX_DATA_HI} swapped & The frame arrives with bytes 4--7 first. Shows
  immediately in a bus analyser, and not at all in simulation if the test is wrong too. \\
The shifts in reverse order & \code{AA BB CC DD} arrives as \code{DD CC BB AA}. \\
\end{booktable}

Both give a frame that \emph{does arrive} --- just with the wrong content. That is the most
unpleasant kind of error, because everything \emph{looks} like it works.

\section{\texorpdfstring{\code{receive()}}{receive()}}
\label{sec:receive}

\begin{cppcode}
bool Spi::receive(Frame& frame)
{
    if ((readReg(RegStatus) & StatusRxValid) == 0U) { return false; }

    frame.id  = static_cast<std::uint16_t>(readReg(RegRxId) & 0x7FFU);
    frame.dlc = static_cast<std::uint8_t>(readReg(RegRxDlc) & 0xFU);

    unpack(readReg(RegRxDataLo), frame.data, 0U);
    unpack(readReg(RegRxDataHi), frame.data, 4U);

    for (std::uint8_t i{frame.dlc}; i < Frame::MaxDataLen; ++i)
    {
        frame.data[i] = 0U;
    }

    writeReg(RegRxAck, Trigger);
    return true;
}
\end{cppcode}

\subsection{The masking against \texorpdfstring{\code{dlc}}{dlc}}

The loop looks redundant. The register map does say that the controller clears its data
accumulator at the start of every frame, so those bytes \emph{are} already zero.

Do it anyway, for three reasons:

\begin{enumerate}
  \item \textbf{The zeroes are the controller's guarantee, not the frame's content.} A driver that
    copies eight bytes straight off reports up to six bytes that were never sent. That they happen
    to be zero does not make them data.
  \item \textbf{The guarantee can disappear.} It comes from an implementation detail in the
    hardware's VHDL. If that changes, or if you borrow a board from another group, your driver is
    still right.
  \item \textbf{\code{dlc} is the contract.} \code{frame.dlc = 2} means "two bytes apply". Code
    that hands out an object whose fields contradict each other leaves the next reader unsure.
\end{enumerate}

\subsection{\texorpdfstring{\code{RX_ACK}}{RX\_ACK} last}

The acknowledgement releases the \code{RX_*} registers and clears \code{STATUS} bit 1. Write it
\textbf{after} you have read everything --- otherwise the next frame can overwrite the registers
in the middle of your reading.

Forget it entirely, and bit 1 stays set forever: every \code{receive()} returns \code{true} with
the same frame, in an endless loop. That is a symptom worth recognising.

\section{\texorpdfstring{\code{hasError()}}{hasError()} and \texorpdfstring{\code{clearError()}}{clearError()}}
\label{sec:errors}

\begin{cppcode}
bool Spi::hasError() const noexcept
{
    return (readReg(RegStatus) & StatusError) != 0U;
}

void Spi::clearError() noexcept
{
    writeReg(RegErrorFlags, 0U);
}
\end{cppcode}

Two lines, and one of them is a trap.

\textbf{\code{clearError()} writes \code{0x0}, not \code{0x1}.} \code{ERROR_FLAGS} is a
latch-and-clear-on-zero construction: only a write where the whole word is zero clears the latch.
\code{0x1} does nothing, and \code{0x000000FF} does nothing even though bit 0 is set.

That is the opposite rule to \code{TX_SEND} and \code{RX_ACK}, which require bit 0 to \emph{be}
set. That the three registers have different rules is not elegant, but it is what the map says,
and it is the map that counts.

\textbf{The symptom if you write \code{0x1}:} the error bit never goes low, \code{hasError()}
returns \code{true} forever after the first error, and all error handling above stops working.

\section{Waiting for a transmission}
\label{sec:waiting}

\code{send()} \textbf{does not wait}. It returns as soon as the trigger is written, and that is
right: a driver method that blocks for up to 130~µs on a 4~MHz MCU eats the whole system's time
budget.

So how does the caller wait? The interface has neither \code{txReady()} nor \code{waitForSend()}
--- deliberately, since \code{Stub} could not have implemented them meaningfully. What the caller
has is that \code{send()} returns \code{false} while the transmitter is busy. \textbf{Waiting is
therefore another attempt:}

\begin{cppcode}
driver.clearError();

if (!isValid(frame)) { return false; }   // never going to be accepted, however long we wait

std::uint16_t attempts{};
while (!driver.send(frame))
{
    if (++attempts >= MaxAttempts) { return false; }
}

// The port was free when the frame was accepted, so the previous transmission had finished.
if (driver.hasError())
{
    driver.clearError();
    // Lost arbitration, stuff error, bad CRC or a remote frame - the latch does not say which.
}
\end{cppcode}

\begin{limitation}
\textbf{A lost arbitration looks like a successful transmission.} \code{STATUS} bit 0 comes back
in both cases, and means \emph{the port is free}, not \emph{the frame arrived}. If you want to
know the outcome you have to read \code{ERROR_FLAGS} \emph{after} the bit has come back --- and
even then you learn \emph{that} something went wrong, never \emph{what}.
\end{limitation}

\section{Checklist}
\label{sec:checklist}

\begin{itemize}
  \item[\checkbox] \code{send()} validates \textbf{first}, and sends no transaction for an invalid
    frame.
  \item[\checkbox] \code{send()} checks \code{StatusTxReady} before it writes anything.
  \item[\checkbox] \code{TX_SEND} is written \textbf{last}.
  \item[\checkbox] \code{TX_DATA_LO} carries data bytes 0--3, with byte 0 in bits 31--24.
  \item[\checkbox] Every \code{static_cast<std::uint32_t>} stands \textbf{before} the shift.
  \item[\checkbox] \code{receive()} checks \code{StatusRxValid} first.
  \item[\checkbox] \code{receive()} writes \code{RX_ACK} \textbf{after} all the reads.
  \item[\checkbox] \code{receive()} zeroes \code{frame.data} above \code{frame.dlc}.
  \item[\checkbox] \code{clearError()} writes \code{0x0}.
  \item[\checkbox] Not a single magic number in the whole file.
\end{itemize}
